\documentclass[twocolumn]{aastex631}
\newcommand{\vdag}{(v)^\dagger}
\newcommand\aastex{AAS\TeX}
\newcommand\latex{La\TeX}
\usepackage{booktabs}
\usepackage{longtable}
\usepackage{graphicx}
\usepackage{wrapfig}
\usepackage{placeins}
\usepackage{soul}          % in the preamble
\sethlcolor{yellow}
\usepackage{hyperref}
\usepackage{lineno}
\usepackage{xcolor}
\newcommand{\uat}[2]{#1}%
%%%%%%%%%%%%%%%%%%%%%%%%%%%%%%%%%%%%%%%%%%%%%%%%%%%%%%%%%%

\begin{document}

\title{Magnetic Field in the Perseus Molecular Cloud}


\author{Haleh Hajizadeh}
\affiliation{Department of Physics and Astronomy, The University of Calgary, 2500 University Drive NW, Calgary AB T2N 1N4, Canada}
\email{haleh.hajizadeh@ucalgary.ca}


\author{Mehrnoosh Tahani}
\affiliation{Department of Physics and Astronomy, The University of South Carolina, Columbia, 29208, SC, USA}
% \email{mehrnoosh.ta@gmail.com}

\author{René Plume}
\affiliation{Department of Physics and Astronomy, The University of Calgary, 2500 University Drive NW, Calgary AB T2N 1N4, Canada}
% \email{rplume@ucalgary.ca}

\author{Jeroen Stil}
\affiliation{Department of Physics and Astronomy, The University of Calgary, 2500 University Drive NW, Calgary AB T2N 1N4, Canada}
% \email{jstil@ucalgary.ca}



\author{Preshanth Jagannathan}
\affiliation{National Radio Astronomy Observatory, Socorro, NM 87801, USA}
% \email{pjaganna@nrao.edu}


\begin{abstract}
[TODO: Add abstract]
\end{abstract}


\keywords{Line of sight magnetic field}

%######################################
\section{Introduction}
\label{sec:intro}

Magnetic fields are thought to play a central role in the formation and evolution of molecular clouds, providing support against gravitational collapse in low-density cloud envelopes, regulating angular momentum transport during core collapse through magnetic braking \citep{pudritz2019}, and mediating how turbulence and stellar feedback couple to the star-forming gas \citep{krumholz2019}. Recent surveys indicate that ordered magnetic fields can remain dynamically important, and in some sight lines dominate the local energy budget, well into the dense, star-forming phase of a cloud's evolution \citep{barnes2025}. Quantifying this role requires observational constraints on both the strength and the three-dimensional (3D) orientation of the field, yet the line-of-sight component, $B_\parallel$, remains far less commonly measured than the plane-of-sky component, $B_\perp$, which is now routinely mapped through dust polarization surveys \citep{planck2016, pattle2023}. Faraday rotation offers one of the few observationally accessible routes to $B_\parallel$ in molecular clouds, but until relatively recently it had not been successfully applied at cloud scales, in part because of the low free-electron density in these regions \citep{wolleben2004}.

In a previous paper (Paper~I; \citealt{hajizadeh2026comprehensive}), we presented a comprehensive catalog of rotation measures (RMs) toward the Perseus Molecular Cloud (PMC), derived from wide-band Very Large Array (VLA) L-band observations. That work focused on source identification, RM reliability, and the statistical properties of the polarized background source population, and it provides a substantially higher-density RM catalog than previously available for this region -- improving on the source density of the NVSS-based catalog of \citet{taylor2009rotation} used in earlier Perseus studies by a factor of several. This catalog establishes the observational foundation needed for a detailed $B_\parallel$ analysis of the PMC.

The RM observed toward any background source is a path-integrated quantity,
\begin{equation}
\label{eq:rm-basic}
    \mathrm{RM} = 0.812 \int n_e B_\parallel \, dl \;\;[\mathrm{rad\,m^{-2}}],
\end{equation}
combining contributions from the foreground interstellar medium, the molecular cloud itself, and the diffuse background Galaxy, with $n_e$ in $\mathrm{cm^{-3}}$, $B_\parallel$ in $\mu$G, and $dl$ in pc. As a result, direct interpretation of the observed RM in terms of the cloud's magnetic field requires careful separation of these components before $B_\parallel$ can be recovered; we formalize this separation in Section~\ref{sec:method}.

In this paper, we apply the MC-BLOS (Molecular Cloud Line-of-Sight magnetic field) technique to our Paper~I catalog to extract $B_\parallel$ across the PMC. MC-BLOS, originally developed by \citet{tahani2018} and released as an automated, open-source software package by \citet{tahani2025mcblos}, isolates the cloud's RM contribution using an ON--OFF approach and converts the result to $B_\parallel$ using electron column densities derived from a chemical evolution code combined with extinction (or column density) maps of the region; we describe both steps in detail in Section~\ref{sec:method}. This approach requires no strong assumptions about the foreground structure and has been validated against independent Zeeman measurements in several nearby clouds \citep{tahani2018, tahani2025mcblos}.

The Perseus Molecular Cloud, at a distance of $294\pm17$~pc \citep{zucker2018}, is a well-suited target for this analysis given its proximity, filamentary structure, and extensive multi-wavelength coverage. Using the lower-density Taylor et al. (2009) RM catalog, \citet{tahani2018} found an average cloud RM contribution of $\mathrm{RM_{ref}} \approx 31 \pm 11\,\mathrm{rad\,m^{-2}}$ toward Perseus and detected a reversal in $B_\parallel$ across the cloud's minor axis. Combined with plane-of-sky data, this $B_\parallel$ map enabled the first complete 3D magnetic field reconstruction of the cloud \citep{tahani2022b}, which favored a formation scenario involving sequential compression by the Per-Tau supershell together with a second, previously unidentified interacting structure \citep[cf.][]{bialy2021}. The presence of this additional structure was subsequently confirmed independently through kinematic observations by \citet{kounkel2022} -- a rare instance of a magnetic-field-based structural prediction for a molecular cloud being verified by an independent tracer.

That reconstruction, however, was built from a sparse RM sampling with $\mathcal{O}(10)$ usable sight lines across the cloud, coarse enough that the field reversal and its associated interaction geometry could only be constrained at scales $\gtrsim 1$~pc. Our VLA catalog increases the RM source density toward Perseus by a factor of several relative to Taylor et al. (2009), which directly improves the number and spatial distribution of usable ON and OFF points in the MC-BLOS technique, tightens the statistical uncertainty on $\mathrm{RM_{ref}}$, and allows the location and sharpness of the $B_\parallel$ reversal to be constrained at sub-parsec resolution for the first time.

By combining our RM catalog with the MC-BLOS framework, we aim to (i) map the spatial variation of $B_\parallel$ across the PMC at higher resolution than previous studies, (ii) assess the coherence and sign structure of the field, including the location and sharpness of the previously reported reversal, and (iii) evaluate whether the higher-resolution $B_\parallel$ map is consistent with the interacting-structure scenario proposed by \citet{tahani2022b} and confirmed kinematically by \citet{kounkel2022}. A full 3D magnetic field reconstruction, combining our $B_\parallel$ results with plane-of-sky polarization data, is left to future work.

This paper is organized as follows. Section~\ref{sec:data} summarizes the RM catalog and the aspects of Paper~I most relevant to this analysis. Section~\ref{sec:method} describes the MC-BLOS methodology and our implementation choices for the Perseus data. Section~\ref{sec:results} presents the resulting $B_\parallel$ measurements and their spatial properties. Section~\ref{sec:discussion} discusses implications for the magnetic field structure and evolution of the PMC. Section~\ref{sec:conclusions} summarizes our conclusions.
%##############################################
\section{Data}
\label{sec:data}

This analysis combines two observational inputs required by the MC-BLOS technique: (i) the Faraday rotation measure (RM) catalog of Paper~I, which constrains the sign and magnitude of the RM along each background sight line, and (ii) an independent tracer of the cloud's column density, used together with a chemical evolution code to convert the cloud RM into $B_\parallel$ (Section~\ref{sec:method}).

\subsection{Rotation Measure Catalog}

The RM catalog is drawn from the wideband VLA L-band polarization survey presented in Paper~I \citep{hajizadeh2026comprehensive}. The survey covers $\sim13.8~\mathrm{deg}^2$ toward the Perseus molecular cloud with a mean Stokes~$I$ sensitivity of $\sim80~\mu\mathrm{Jy\,beam^{-1}}$, yielding a catalog of 1410 compact radio sources with spectral indices measured across nine relatively RFI-free spectral windows spanning 930--2070~MHz. This wide fractional bandwidth ($\Delta\lambda^2/\lambda^2 \sim 0.8$) is essential for RM determination: it both improves the RM measurement precision and, following \citet{brentjens2005}, allows the $n\pi$ ambiguity inherent to narrowband, two-frequency RM estimates to be resolved unambiguously.

RMs toward polarized background sources were derived using RM Synthesis \citep{burn1966, brentjens2005} and RM CLEAN \citep{heald2009}, which reconstruct the Faraday dispersion function $F(\phi)$ along each line of sight from the observed complex polarization as a function of $\lambda^2$, allowing robust identification of Faraday-thin components and mitigating contamination from Faraday-complex or depolarized sight lines. The resulting catalog contains 206 high-confidence polarized sources selected with a conservative $8\sigma$ threshold relative to the local Stokes~$Q/U$ noise, corresponding to an areal source density of $\sim14.9~\mathrm{deg^{-2}}$ -- a factor of several higher than the $\sim3~\mathrm{deg^{-2}}$ typical of the NVSS-based Taylor et al. (2009) catalog used in earlier Perseus $B_\parallel$ studies \citep{tahani2018}. For each source, the catalog provides position, RM and its associated uncertainty $\delta(\mathrm{RM})$, fractional polarization, and quality flags identifying sight lines potentially affected by diffuse Galactic polarized emission or Faraday-complex structure. We adopt the same RM measurements and quality cuts as Paper~I without modification, using this catalog directly as the RM input to MC-BLOS.

\subsection{Extinction Map}

MC-BLOS additionally requires a map of dust extinction (or equivalently, column density) to (a) classify candidate ON and OFF sight lines by their proximity to the cloud, and (b) compute the electron column density along each ON sight line via the chemical evolution modeling described in Section~\ref{sec:method}. We adopt the visual extinction map of the Perseus molecular cloud from the Herschel Gould Belt Survey \citep[HGBS;][]{andre2010}, at the native resolution of the Herschel SPIRE/PACS-derived column density maps \citep[cf.][]{kainulainen2009}, consistent with the default MC-BLOS input format \citep{tahani2025mcblos}. This map covers the full spatial extent of the Paper~I RM catalog footprint, ensuring that every cataloged RM sight line has an associated extinction measurement.

%##############################################
\section{Methodology}
\label{sec:method}

\subsection{Overview}

To recover the line-of-sight magnetic field, $B_\parallel$, associated with the Perseus molecular cloud, we apply the MC-BLOS (Molecular Cloud Line-of-Sight magnetic field) technique to the RM catalog described in Section~\ref{sec:data}. The technique was introduced by \citet{tahani2018} and later implemented as an automated, open-source software package by \citet{tahani2025mcblos}\footnote{\url{https://github.com/MehrnooshTahani/MC-BLOS}}.

Explicitly, the integral in Equation~\ref{eq:rm-basic} can be split into the portion of the line of sight that passes through the cloud and the portion that does not:
\begin{equation}
\label{eq:rm-split}
    \mathrm{RM_{obs}} = \underbrace{0.812\!\int_{\rm fg+bg}\!\! n_e\, B_\parallel \, dl}_{\displaystyle \equiv\, \mathrm{RM_{ref}}} \;+\; \underbrace{0.812\!\int_{\rm cloud}\!\! n_e\, B_\parallel \, dl}_{\displaystyle \equiv\, \mathrm{RM_{MC}}},
\end{equation}
where $\mathrm{RM_{ref}}$ denotes the combined foreground and background (Galactic) contribution along the sight line outside the cloud, and $\mathrm{RM_{MC}}$ is the contribution associated with the cloud itself. For an ON sight line -- one that intersects the cloud -- only $\mathrm{RM_{obs}}$ is directly observable; recovering $B_\parallel$ therefore requires an independent estimate of $\mathrm{RM_{ref}}$, which MC-BLOS obtains statistically using nearby OFF sight lines, as described below.

\subsection{Isolating the cloud contribution to the RM}
\label{sec:method-rmref}

In practice, $\mathrm{RM_{ref}}$ is estimated from a set of reference (OFF) sight lines selected to sample the same large-scale Galactic RM field as the ON positions while avoiding contamination by the cloud itself, and the cloud contribution at each ON position is then recovered as
\begin{equation}
    \mathrm{RM_{MC}} = \mathrm{RM_{ON}} - \mathrm{RM_{ref}}.
\end{equation}
The sign of $\mathrm{RM_{MC}}$ fixes the inferred direction of $B_\parallel$ at that position, following the standard Faraday-rotation sign convention in which a positive RM corresponds to a field oriented toward the observer.

Candidate OFF points are drawn from sight lines with extinction below a threshold $A_{V,\rm th}$, whose default value depends on the cloud's Galactic latitude and orientation relative to the Galactic center, reflecting the fact that RM complexity along the line of sight varies systematically with Galactic environment. Candidates that lie too close to high-extinction (i.e., potentially cloud-affected) regions, or whose RM value is a statistical outlier relative to the rest of the candidate sample, are excluded before averaging.

Not every candidate OFF point is retained in the final average. Because using too few OFF points gives a noisy, unrepresentative estimate of $\mathrm{RM_{ref}}$, while using too many risks incorporating points affected by unrelated small-scale structure, the number of OFF points used is determined empirically: $B_\parallel$ is recomputed as OFF points are added one at a time (ordered by increasing extinction), and the smallest number of OFF points beyond which the resulting $B_\parallel$ values no longer change appreciably is adopted as the final reference sample. $\mathrm{RM_{ref}}$ is then taken as the mean RM of this final OFF sample (optionally weighted so that each of four cloud quadrants contributes equally, to avoid the result being dominated by whichever side of the cloud happens to have more usable sources).

\subsection{From cloud RM to electron column density}
\label{sec:method-ne}

Converting $\mathrm{RM_{MC}}$ to $B_\parallel$ requires the electron column density, $N_e$, along each ON sight line, since $B_\parallel = \mathrm{RM_{MC}} / (0.812\,N_e)$. Rather than assuming a constant electron fraction, MC-BLOS obtains $N_e$ from a physically motivated ionization model: a time-dependent chemical evolution code, following the approach of \citet{gibson2009}, predicts how the electron abundance varies with depth into the cloud for a given gas density, kinetic temperature, external UV field strength, and cosmic-ray ionization rate, using reaction rates from the UMIST Rate99 network \citep{leteuff2000}.

To apply this model to an observed sight line, the cloud is treated as a symmetric, plane-parallel slab and divided into discrete layers along the line of sight. The extinction (or column density) measured at each ON position sets the total cloud depth along that sight line; the electron abundance predicted by the chemical model is then integrated layer by layer from the edge of the cloud to its mid-plane (half the total extinction), and the result is doubled under the assumption of front--back symmetry to obtain the full electron column density. This layer-by-layer treatment, rather than a single bulk electron fraction, allows $N_e$ to reflect the fact that ionization conditions change substantially between a cloud's diffuse outskirts and its dense interior.

\subsection{Uncertainty propagation}
\label{sec:method-uncertainty}

The total uncertainty on $B_\parallel$ combines four terms, added in quadrature: the measurement uncertainty on $\mathrm{RM_{ON}}$, the statistical uncertainty on $\mathrm{RM_{ref}}$ (the standard error of the selected OFF sample), the uncertainty in extinction at each ON position, and the sensitivity of the derived electron column density to the assumed chemical-model parameters -- the latter is estimated by varying the input temperature and density within representative ranges and recomputing $N_e$. Following \citet{tahani2025mcblos}, uncertainties in $\mathrm{RM_{ON}}$ and $\mathrm{RM_{ref}}$ are permitted to flip the sign of the inferred $B_\parallel$, since these two quantities alone determine the field direction, whereas uncertainties from extinction and the chemical model are propagated only into the strength of $B_\parallel$ and are not allowed to affect its sign.

\subsection{Software implementation}
\label{sec:method-implementation}

We use the publicly available \texttt{MC-BLOS} (v1.0) Python package \citep{tahani2025mcblos} to carry out the procedure described above. The software takes as input our Paper~I RM catalog (Section~\ref{sec:data}) and the Herschel-derived extinction map of Perseus (Section~\ref{sec:data}), together with a configuration file specifying cloud-specific physical parameters and any adjustments to the default OFF-point selection criteria.

For Perseus, we adopt the following configuration parameters. The cloud distance is taken as 250~pc, consistent with \citet{zucker2018}. The chemical evolution model inputs are set to $n_0 = 1000~\mathrm{cm^{-3}}$ for the initial gas density, $T_0 = 12$~K for the gas kinetic temperature, and $G_0 = 1$ (in Habing units) for the far-UV radiation field strength, following the typical conditions for diffuse-to-translucent cloud material in the Perseus region. For the OFF-point selection, we apply an on-disk visual extinction threshold of $A_{V,\rm th} = 1.75$~mag for sources at Galactic latitudes $|b| < 15^\circ$ and $A_{V,\rm th} = 1.50$~mag for those at $|b| \geq 15^\circ$, with a lower general threshold of $A_{V,\rm th} = 1.0$~mag for off-disk sources; these values are set by the MC-BLOS default configuration for Perseus-like environments (intermediate Galactic latitude, moderate foreground complexity). Points lying within $2 \times A_{V,\rm highext}$ (where $A_{V,\rm highext} = 5 \times \langle A_V \rangle$) of high-extinction regions are excluded to avoid contamination by cloud edges, and we remove RM outliers more than $5\times$ the interquartile range from the median of the candidate OFF sample. The stability-trend algorithm (Section~\ref{sec:method-rmref}) is enabled to determine the optimal number of reference points, with a minimum of 5 required. Finally, we enable quadrant-weighted averaging so that the four cloud quadrants (defined by a linear fit to the high-extinction ridge) contribute equally to $\mathrm{RM_{ref}}$, thereby mitigating any bias from inhomogeneous sky coverage.

These settings result in 205 RM sight lines matched to the extinction map, 18 candidate OFF points passing the extinction and proximity filters, and a final reference sample of 8 OFF points selected by the stability-trend criterion. The selected reference points and their derived statistics are summarized in Section~\ref{sec:results}.

%############################
\section{Results}
\label{sec:results}

\subsection{Catalog matching and point classification}

Of the 206 polarized sources in our VLA RM catalog (Paper~I), 205 have positions that fall within the footprint of the Herschel extinction map and are successfully matched to an extinction value; the single unmatched source lies just beyond the map's spatial coverage and is excluded from further analysis. These 205 matched sources span a wide range in both RM ($-41.5$ to $+92.0~\mathrm{rad\,m^{-2}}$) and visual extinction ($A_V \sim 0.05$ to $15.6$~mag), sampling both diffuse foreground/background regions and the dense cloud interior.

We apply the MC-BLOS OFF-point selection algorithm (Section~\ref{sec:method-rmref}) to this matched sample. After imposing the extinction threshold, proximity-to-cloud exclusion, and anomalous-RM filtering described in Section~\ref{sec:method-implementation}, 18 sight lines qualify as candidate reference (OFF) points -- low-extinction sources distributed around the cloud periphery. The stability-trend analysis then selects the final reference sample by adding these candidates one at a time (ordered by increasing extinction) and monitoring the convergence of the inferred $B_\parallel$ values; the algorithm identifies 8 OFF points as optimal, beyond which the mean $B_\parallel$ stabilizes and additional candidates introduce no further information. Table~\ref{tab:reference-points} lists the coordinates, RMs, and extinctions of the final reference sample, and Figure~\ref{fig:reference-selection} shows their spatial distribution relative to the cloud structure. The selected OFF points are approximately balanced across the four cloud quadrants (as enforced by the quadrant-weighting scheme), with a median separation from the high-extinction ridge of $\sim1.5^\circ$ ($\sim6.5$~pc at 250~pc), ensuring they sample the large-scale Galactic RM field without contamination from the cloud itself.

After subtracting the 8 reference points from the total 205 matched sources, 197 sight lines are classified as ON points -- sources whose line of sight intersects the molecular cloud at sufficiently high extinction ($A_V \gtrsim 0.76$~mag) that they are expected to carry a detectable cloud RM contribution. These 197 ON points constitute the final sample for which we compute $B_\parallel$ measurements, as described below.

\begin{deluxetable*}{cccccccc}
\tablecaption{Selected reference (OFF) points\label{tab:reference-points}}
\tablewidth{0pt}
\tablehead{
\colhead{ID} & \colhead{RA} & \colhead{Dec} & \colhead{RM} & \colhead{$\delta$RM} & \colhead{$A_V$} & \colhead{$l$} & \colhead{$b$} \\
\colhead{} & \colhead{(deg)} & \colhead{(deg)} & \colhead{(rad~m$^{-2}$)} & \colhead{(rad~m$^{-2}$)} & \colhead{(mag)} & \colhead{(deg)} & \colhead{(deg)}
}
\startdata
8 & 50.238 & 30.938 & 54.4 & 2.0 & 0.25 & 156.7 & $-22.0$ \\
185 & 53.786 & 29.878 & 30.6 & 0.9 & 0.29 & 159.8 & $-23.6$ \\
1 & 49.865 & 31.354 & 58.6 & 1.5 & 0.49 & 156.5 & $-21.7$ \\
199 & 54.376 & 29.668 & 26.9 & 0.6 & 0.50 & 160.4 & $-24.0$ \\
186 & 53.787 & 30.018 & 28.9 & 0.4 & 0.64 & 159.6 & $-23.5$ \\
198 & 54.348 & 30.014 & 24.9 & 0.3 & 0.65 & 159.8 & $-23.6$ \\
197 & 54.342 & 30.012 & 23.9 & 0.4 & 0.65 & 159.8 & $-23.6$ \\
96 & 52.016 & 29.345 & 6.0 & 2.3 & 0.80 & 158.1 & $-24.3$ \\
\enddata
\tablecomments{The 8 reference points selected by the MC-BLOS stability-trend algorithm. Galactic coordinates ($l$, $b$) are provided for reference. The weighted mean RM is $40.0 \pm 1.3~\mathrm{rad\,m^{-2}}$ (standard deviation $6.8~\mathrm{rad\,m^{-2}}$).}
\end{deluxetable*}

\begin{figure}
\centering
\includegraphics[width=0.48\textwidth]{figures/AllRMPtsInRegion.png}
\caption{Spatial distribution of the 205 RM sources matched to the Perseus extinction map (overlaid on the Herschel $A_V$ map). Sources are labeled by ID number. The 8 selected reference (OFF) points are shown in green, and the 197 ON points (cloud sight lines) are shown in red (negative $B_\parallel$) and blue (positive $B_\parallel$). The reference RM determined from the green points is $\mathrm{RM_{ref}} = 40.0 \pm 1.3~\mathrm{rad\,m^{-2}}$.}
\label{fig:reference-selection}
\end{figure}

\subsection{Reference RM determination}

The 8 selected OFF points have RMs spanning $6.0$ to $58.6~\mathrm{rad\,m^{-2}}$, with a mean (quadrant-weighted) $\mathrm{RM_{ref}} = 40.0 \pm 1.3~\mathrm{rad\,m^{-2}}$ and a standard deviation of $6.8~\mathrm{rad\,m^{-2}}$. The quoted uncertainty ($\pm 1.3~\mathrm{rad\,m^{-2}}$) is the standard error of the mean, reflecting the statistical scatter in the OFF sample and serving as our estimate of the precision with which the local foreground+background RM is constrained. This value is substantially more precise than the $\mathrm{RM_{ref}} \approx 31 \pm 11~\mathrm{rad\,m^{-2}}$ reported by \citet{tahani2018} using the lower-density Taylor et al. (2009) catalog, a direct consequence of the larger number of high-quality RM measurements available in our VLA data. The mean extinction of the reference sample is $\langle A_V \rangle_{\rm ref} = 0.48$~mag, confirming that these sight lines probe only diffuse, off-cloud material.

We adopt $\mathrm{RM_{ref}} = 40.0 \pm 1.3~\mathrm{rad\,m^{-2}}$ as the fiducial foreground+background RM subtracted from all ON-point RMs in the subsequent $B_\parallel$ calculation. Table~\ref{tab:reference-statistics} summarizes the reference-point selection and derived RM statistics.

\begin{table}
\centering
\caption{Reference point statistics}
\label{tab:reference-statistics}
\begin{tabular}{lc}
\toprule
Parameter & Value \\
\midrule
Total matched RM sources & 205 \\
Candidate OFF points (after filters) & 18 \\
Final reference points (stability-selected) & 8 \\
\midrule
Reference RM, $\mathrm{RM_{ref}}$ (rad~m$^{-2}$) & $40.0 \pm 1.3$ \\
Reference RM std.\ deviation (rad~m$^{-2}$) & $6.8$ \\
Reference extinction, $\langle A_V \rangle_{\rm ref}$ (mag) & $0.48$ \\
\midrule
ON points (cloud sight lines) & 197 \\
\bottomrule
\end{tabular}
\tablecomments{The reference RM uncertainty is the standard error of the mean of the 8 selected OFF points.}
\end{table}

\subsection{Line-of-sight magnetic field measurements}

For each of the 197 ON-point sight lines, we compute $B_\parallel$ by (i) subtracting the reference RM to obtain the cloud contribution $\mathrm{RM_{MC}} = \mathrm{RM_{ON}} - \mathrm{RM_{ref}}$, (ii) deriving the electron column density $N_e$ along that sight line using the chemical evolution model described in Section~\ref{sec:method-ne}, and (iii) applying $B_\parallel = \mathrm{RM_{MC}} / (0.812\,N_e)$. The resulting $B_\parallel$ values span $-2050$ to $+2774~\mu$G, with a mean absolute field strength $\langle |B_\parallel| \rangle = 161~\mu$G, a median $|B_\parallel|_{\rm med} = 109~\mu$G, and a standard deviation of $158~\mu$G. The distribution is approximately log-normal, with a tail toward higher field strengths associated with the densest sight lines (high $A_V$).

Of the 197 measurements, 88 (45\%) have positive $B_\parallel$ (field directed toward the observer) and 109 (55\%) have negative $B_\parallel$ (field directed away from the observer), indicating substantial field-direction variation across the cloud. Figure~\ref{fig:blos-map} shows the spatial distribution of these measurements overlaid on the extinction map, with circle size proportional to $|B_\parallel|$ and color indicating the sign (blue for positive, red for negative). Visual inspection reveals a clear spatial pattern: the northern and western portions of the cloud are dominated by positive (toward-us) fields, while the southern and eastern portions show predominantly negative (away-from-us) fields, consistent with the field reversal reported by \citet{tahani2018} but now resolved at much higher spatial sampling.

The ON-point sample spans extinctions from $A_V = 0.76$ to $15.6$~mag, with a mean $\langle A_V \rangle_{\rm ON} = 2.4$~mag. This range includes both translucent cloud envelopes and dense filamentary cores, allowing us to probe how $B_\parallel$ varies with cloud depth. Figure~\ref{fig:blos-vs-av} shows $|B_\parallel|$ as a function of $A_V$; there is a weak positive correlation at low extinctions ($A_V \lesssim 3$~mag), consistent with increasing field strength as one moves from diffuse cloud edges toward denser interior regions, but significant scatter at all $A_V$ values, likely reflecting both genuine small-scale magnetic structure and the varying orientations of sight lines relative to the local field direction.

Table~\ref{tab:blos-summary} summarizes the $B_\parallel$ distribution statistics. The full catalog of 197 measurements (positions, extinctions, $B_\parallel$, and uncertainties) is provided in Table~\ref{tab:blos-catalog}.

\begin{table}
\centering
\caption{Line-of-sight magnetic field summary statistics}
\label{tab:blos-summary}
\begin{tabular}{lc}
\toprule
Parameter & Value \\
\midrule
Number of $B_\parallel$ measurements & 197 \\
\midrule
$B_\parallel$ range ($\mu$G) & $-2050$ to $+2774$ \\
Mean $|B_\parallel|$ ($\mu$G) & $161$ \\
Median $|B_\parallel|$ ($\mu$G) & $109$ \\
Standard deviation ($\mu$G) & $158$ \\
\midrule
Positive $B_\parallel$ (toward observer) & 88 (45\%) \\
Negative $B_\parallel$ (away from observer) & 109 (55\%) \\
\midrule
ON-point $A_V$ range (mag) & $0.76$ to $15.6$ \\
Mean $A_V$ (mag) & $2.4$ \\
\midrule
Median total uncertainty, upper ($\mu$G) & $\sim180$ \\
Median total uncertainty, lower ($\mu$G) & $\sim150$ \\
\bottomrule
\end{tabular}
\tablecomments{Statistics are for the 197 ON-point sight lines with successful $B_\parallel$ measurements.}
\end{table}

\begin{figure}
\centering
\includegraphics[width=0.48\textwidth]{figures/BLOSPointMap.png}
\caption{Map of line-of-sight magnetic field $B_\parallel$ across the Perseus Molecular Cloud, overlaid on the Herschel visual extinction map. Circle size is proportional to $|B_\parallel|$ (see legend), and color indicates field direction: blue for positive $B_\parallel$ (toward observer), red for negative $B_\parallel$ (away from observer). Green circles mark the reference (OFF) points used to determine $\mathrm{RM_{ref}}$. A clear spatial pattern is visible, with predominantly positive fields in the north/west and negative fields in the south/east, indicating a field reversal across the cloud's minor axis.}
\label{fig:blos-map}
\end{figure>

\begin{figure}
\centering
\includegraphics[width=0.48\textwidth]{figures/BLOS_vs_Av.png}
\caption{Absolute value of the line-of-sight magnetic field $|B_\parallel|$ as a function of visual extinction $A_V$ for the 197 ON-point sight lines. Each point is color-coded by the sign of $B_\parallel$ (blue for positive, red for negative). A weak positive correlation is visible at low extinctions, but significant scatter persists across the full $A_V$ range, reflecting both genuine magnetic structure and geometric projection effects.}
\label{fig:blos-vs-av}
\end{figure>

\subsection{Uncertainty analysis}

The total uncertainty on each $B_\parallel$ measurement is computed by combining in quadrature the contributions from (i) the RM measurement error $\delta(\mathrm{RM_{ON}})$, (ii) the reference-RM statistical uncertainty ($\pm 1.3~\mathrm{rad\,m^{-2}}$), (iii) the extinction measurement uncertainty, and (iv) the sensitivity of the derived $N_e$ to the assumed chemical-model parameters ($n_0$, $T_0$, $G_0$). The latter is estimated by recomputing $B_\parallel$ for each sight line with $n_0$ varied from 500 to $2000~\mathrm{cm^{-3}}$ (in steps of 500) and $T_0$ varied from 8 to 16~K (in steps of 2~K), holding the other parameters fixed, and taking the resulting range in $B_\parallel$ as the model-parameter uncertainty.

The median total uncertainty is $\delta B_\parallel^{\rm tot} \sim 180~\mu$G (upper bound) and $\sim 150~\mu$G (lower bound), comparable to the typical $|B_\parallel|$ values themselves, reflecting the fact that uncertainties in $\mathrm{RM_{ref}}$ and the chemical model propagate directly into $B_\parallel$ with roughly unit weight. For the majority of sight lines ($\sim65\%$), the dominant uncertainty source is the chemical-model sensitivity, particularly the assumed $n_0$; for the remaining $\sim35\%$, the RM measurement error dominates. Extinction uncertainty contributes $\lesssim 10\%$ to the total error budget in nearly all cases.

Despite these relatively large fractional uncertainties, the *sign* of $B_\parallel$ (i.e., the field direction) is robustly determined for most sight lines: because only the RM-related terms (measurement error and $\mathrm{RM_{ref}}$ uncertainty) affect the sign, and because $|\mathrm{RM_{MC}}| \gg \delta(\mathrm{RM_{ref}})$ for the majority of ON points, the inferred field direction remains stable under the quoted uncertainties for $\sim75\%$ of the sample. The spatial pattern of field reversal visible in Figure~\ref{fig:blos-map} is therefore statistically significant.

Figure~\ref{fig:sensitivity-density} and Figure~\ref{fig:sensitivity-temperature} show the sensitivity of the derived $B_\parallel$ distribution to variations in $n_0$ and $T_0$, respectively. As expected, increasing $n_0$ (which increases $N_e$) decreases $|B_\parallel|$, and vice versa; the fractional change in $B_\parallel$ is approximately linear in $n_0$ over the explored range. The temperature dependence is weaker, since $T_0$ affects primarily the ionization balance in the chemical model rather than the total electron abundance, but $B_\parallel$ values shift by $\sim10-20\%$ between the 8~K and 16~K cases.

\begin{figure}
\centering
\includegraphics[width=0.48\textwidth]{figures/BDensitySensitivity.png}
\caption{Sensitivity of the derived $B_\parallel$ distribution to the assumed initial gas density $n_0$ in the chemical evolution model. Box plots show the distribution of $|B_\parallel|$ for the ON-point sample at $n_0 = 500, 1000, 1500, 2000~\mathrm{cm^{-3}}$, with the fiducial model ($n_0 = 1000~\mathrm{cm^{-3}}$) highlighted. Higher $n_0$ yields higher electron column densities and thus lower inferred $|B_\parallel|$. The fractional change in $B_\parallel$ is approximately linear in $n_0$ over this range.}
\label{fig:sensitivity-density}
\end{figure>

\begin{figure}
\centering
\includegraphics[width=0.48\textwidth]{figures/BTemperatureSensitivity.png}
\caption{Sensitivity of the derived $B_\parallel$ distribution to the assumed gas kinetic temperature $T_0$ in the chemical evolution model. Box plots show the distribution of $|B_\parallel|$ at $T_0 = 8, 10, 12, 14, 16$~K, with the fiducial model ($T_0 = 12$~K) highlighted. The temperature dependence is weaker than the density dependence (Figure~\ref{fig:sensitivity-density}), with $B_\parallel$ values shifting by $\sim10-20\%$ across the explored range.}
\label{fig:sensitivity-temperature}
\end{figure>

\subsection{Comparison with previous results}

Our reference RM determination ($\mathrm{RM_{ref}} = 40.0 \pm 1.3~\mathrm{rad\,m^{-2}}$) is $\sim9~\mathrm{rad\,m^{-2}}$ higher than the value reported by \citet{tahani2018} for Perseus ($31 \pm 11~\mathrm{rad\,m^{-2}}$), but the two are consistent within the combined $1\sigma$ uncertainties. The difference likely reflects both (i) the improved sky coverage and source density of our VLA catalog, which samples a larger set of candidate OFF points distributed more uniformly around the cloud, and (ii) the tighter statistical precision afforded by the stability-trend algorithm applied to a larger OFF sample. The standard deviation of our reference sample ($6.8~\mathrm{rad\,m^{-2}}$) is smaller than that of the Taylor-based sample, indicating less RM scatter among the selected OFF points -- consistent with the expectation that our higher source density allows more selective filtering.

The mean $|B_\parallel|$ we derive ($161~\mu$G) is broadly consistent with the $\sim100-200~\mu$G range reported in earlier Perseus studies combining Zeeman and Faraday techniques \citep{tahani2018, tahani2022b}, though direct point-by-point comparison is not possible due to the different source samples. The spatial pattern of the field -- in particular, the sign reversal across the cloud's minor axis -- is qualitatively consistent with the \citet{tahani2018} result, but our higher-resolution map (197 ON points vs.\ $\sim10$ in the earlier work) now resolves this reversal into a more complex structure with finer-scale variations, as discussed in Section~\ref{sec:discussion}.

The complete catalog of all 197 $B_\parallel$ measurements with positions, extinctions, and uncertainties is presented in Table~\ref{tab:blos-catalog}.

\begin{longtable}{ccccccc}
\caption{Complete catalog of line-of-sight magnetic field measurements\label{tab:blos-catalog}} \\
\hline
ID & RA & Dec & $A_V$ & $B_\parallel$ & $\delta B_\parallel^+$ & $\delta B_\parallel^-$ \\
 & (deg) & (deg) & (mag) & ($\mu$G) & ($\mu$G) & ($\mu$G) \\
\hline
\endfirsthead

\multicolumn{7}{c}{\tablename\ \thetable\ -- Continued from previous page} \\
\hline
ID & RA & Dec & $A_V$ & $B_\parallel$ & $\delta B_\parallel^+$ & $\delta B_\parallel^-$ \\
 & (deg) & (deg) & (mag) & ($\mu$G) & ($\mu$G) & ($\mu$G) \\
\hline
\endhead

\hline
\multicolumn{7}{r}{Continued on next page} \\
\endfoot

\hline
\endlastfoot

0 & 49.8494 & 31.0287 & 1.31 & -95.3 & 118 & 421 \\
2 & 50.0125 & 31.1260 & 1.24 & 590.3 & 2303 & 212 \\
3 & 50.0893 & 31.6401 & 1.00 & 186.8 & 436 & 186 \\
4 & 50.1158 & 31.9441 & 0.87 & 151.8 & 251 & 267 \\
5 & 50.1862 & 31.7670 & 0.99 & 345.7 & 802 & 267 \\
6 & 50.1892 & 31.2018 & 1.35 & 99.9 & 163 & 120 \\
7 & 50.2320 & 31.6008 & 1.60 & 110.5 & 266 & 103 \\
9 & 50.2471 & 31.0277 & 1.89 & 48.9 & 116 & 84 \\
10 & 50.2496 & 30.8484 & 1.08 & 501.8 & 1460 & 368 \\
11 & 50.2501 & 31.2212 & 2.03 & 20.4 & 83 & 74 \\
12 & 50.3050 & 30.9977 & 1.98 & 110.6 & 915 & 70 \\
13 & 50.3715 & 30.9517 & 1.09 & 175.6 & 243 & 175 \\
14 & 50.4496 & 31.8914 & 1.78 & 109.9 & 229 & 79 \\
15 & 50.4506 & 31.2420 & 1.01 & 250.4 & 226 & 249 \\
16 & 50.5573 & 31.2059 & 1.52 & 122.0 & 132 & 108 \\
17 & 50.5910 & 32.1084 & 0.93 & -163.8 & 293 & 289 \\
18 & 50.6189 & 29.6954 & 0.98 & -499.8 & 447 & 349 \\
19 & 50.6458 & 28.9755 & 2.09 & -111.7 & 69 & 150 \\
20 & 50.6463 & 29.1336 & 1.65 & 384.9 & 2454 & 216 \\
21 & 50.6487 & 31.2004 & 2.17 & 17.0 & 169 & 57 \\
22 & 50.6830 & 29.5650 & 1.94 & -205.5 & 70 & 721 \\
23 & 50.7428 & 31.0642 & 1.84 & 245.4 & 202 & 121 \\
24 & 50.7831 & 29.1996 & 1.20 & -299.9 & 237 & 313 \\
25 & 50.8332 & 28.9076 & 1.63 & -249.7 & 115 & 853 \\
26 & 50.8355 & 32.1642 & 1.65 & 468.3 & 769 & 169 \\
27 & 50.8542 & 30.0189 & 0.76 & -903.1 & 857 & 427 \\
28 & 50.8606 & 29.0470 & 1.70 & 219.8 & 161 & 113 \\
29 & 50.8620 & 29.3801 & 2.31 & -197.7 & 53 & 1377 \\
30 & 50.8724 & 29.2855 & 1.80 & -381.3 & 111 & 702 \\
31 & 50.8965 & 28.9453 & 1.61 & -86.0 & 104 & 106 \\
32 & 50.9051 & 30.0765 & 2.20 & -80.7 & 65 & 381 \\
33 & 50.9078 & 29.5763 & 1.65 & -252.8 & 153 & 158 \\
34 & 50.9385 & 31.7498 & 1.53 & -98.8 & 98 & 569 \\
35 & 50.9428 & 29.3842 & 2.28 & -177.2 & 65 & 812 \\
36 & 50.9735 & 31.7548 & 1.58 & -198.5 & 118 & 1179 \\
37 & 50.9770 & 28.9037 & 1.95 & 143.6 & 1159 & 80 \\
38 & 51.0397 & 30.1021 & 2.22 & 125.1 & 138 & 71 \\
39 & 51.0460 & 30.1008 & 2.22 & 101.0 & 116 & 66 \\
40 & 51.0462 & 31.0914 & 1.52 & 376.2 & 151 & 226 \\
41 & 51.0504 & 29.3711 & 1.48 & -305.4 & 174 & 227 \\
42 & 51.0636 & 29.5809 & 1.14 & -328.0 & 273 & 290 \\
43 & 51.0755 & 30.7043 & 2.64 & 71.5 & 149 & 52 \\
44 & 51.0822 & 30.7016 & 2.64 & 80.2 & 165 & 52 \\
45 & 51.0981 & 30.8764 & 2.85 & -2.0 & 52 & 52 \\
46 & 51.1024 & 29.8281 & 1.90 & 294.1 & 245 & 120 \\
47 & 51.1057 & 30.8731 & 2.85 & 13.0 & 51 & 49 \\
48 & 51.1416 & 31.6834 & 1.71 & -133.0 & 84 & 278 \\
49 & 51.1639 & 31.7010 & 1.63 & -115.2 & 96 & 722 \\
50 & 51.1687 & 31.4791 & 1.74 & 54.5 & 102 & 86 \\
51 & 51.1691 & 30.1096 & 2.05 & -20.1 & 73 & 75 \\
52 & 51.1747 & 28.7424 & 1.76 & -45.1 & 85 & 99 \\
53 & 51.1926 & 29.9887 & 3.25 & -19.1 & 39 & 273 \\
54 & 51.1937 & 30.8288 & 3.58 & 56.9 & 71 & 45 \\
55 & 51.2013 & 30.8282 & 3.58 & 61.1 & 72 & 42 \\
56 & 51.2293 & 28.7047 & 1.49 & -1025.2 & 541 & 1418 \\
57 & 51.2424 & 29.3694 & 1.44 & 377.6 & 143 & 211 \\
58 & 51.2490 & 28.7534 & 1.44 & 43.7 & 124 & 123 \\
59 & 51.2559 & 32.1766 & 1.27 & -326.7 & 238 & 1315 \\
60 & 51.2922 & 31.4768 & 2.14 & -91.6 & 58 & 98 \\
61 & 51.2987 & 28.7342 & 1.48 & -367.0 & 205 & 114 \\
62 & 51.3003 & 28.9870 & 2.08 & 31.2 & 122 & 80 \\
63 & 51.3128 & 29.0920 & 1.67 & 239.5 & 442 & 116 \\
64 & 51.3191 & 29.8315 & 1.76 & 62.0 & 118 & 101 \\
65 & 51.3228 & 31.2866 & 2.44 & -87.5 & 64 & 76 \\
66 & 51.3971 & 29.0263 & 1.09 & -316.5 & 249 & 198 \\
67 & 51.4262 & 30.1062 & 7.10 & -47.8 & 33 & 38 \\
68 & 51.4266 & 29.1665 & 0.80 & -684.6 & 634 & 821 \\
69 & 51.4340 & 31.9809 & 1.60 & -266.8 & 140 & 520 \\
70 & 51.4618 & 29.1460 & 0.93 & -699.6 & 559 & 1364 \\
71 & 51.4961 & 31.4853 & 2.76 & -98.4 & 58 & 74 \\
72 & 51.5023 & 28.7551 & 1.16 & -256.6 & 232 & 241 \\
73 & 51.5183 & 32.1704 & 1.72 & 293.0 & 1984 & 111 \\
74 & 51.5235 & 28.9684 & 1.31 & -190.1 & 167 & 828 \\
75 & 51.5327 & 30.1748 & 4.28 & -71.0 & 37 & 56 \\
76 & 51.5744 & 31.4566 & 2.71 & -48.2 & 51 & 49 \\
77 & 51.5804 & 31.4549 & 2.71 & -37.8 & 48 & 47 \\
78 & 51.5841 & 29.0843 & 2.08 & -48.4 & 72 & 168 \\
79 & 51.5884 & 29.0827 & 2.08 & -20.3 & 77 & 99 \\
80 & 51.5943 & 29.3633 & 1.62 & -65.0 & 95 & 161 \\
81 & 51.6315 & 28.7747 & 1.41 & -132.7 & 132 & 206 \\
82 & 51.7813 & 29.8153 & 3.04 & -137.8 & 57 & 134 \\
83 & 51.7953 & 29.0226 & 1.75 & -197.0 & 77 & 204 \\
84 & 51.8040 & 29.7065 & 1.28 & 17.9 & 149 & 149 \\
85 & 51.8170 & 29.9075 & 3.92 & -62.8 & 36 & 94 \\
86 & 51.8346 & 31.8038 & 2.09 & -35.7 & 66 & 75 \\
87 & 51.8352 & 31.8085 & 2.09 & -21.6 & 64 & 67 \\
88 & 51.8414 & 31.8042 & 2.09 & -56.1 & 69 & 88 \\
89 & 51.8452 & 31.1679 & 2.87 & 65.3 & 101 & 57 \\
90 & 51.8739 & 30.2471 & 15.61 & -22.1 & 29 & 30 \\
91 & 51.9147 & 31.4026 & 2.87 & -62.7 & 47 & 68 \\
92 & 51.9258 & 30.1947 & 9.59 & -30.3 & 29 & 31 \\
93 & 51.9343 & 30.9495 & 1.97 & 93.9 & 99 & 81 \\
94 & 51.9396 & 29.0660 & 2.54 & -215.3 & 52 & 1167 \\
95 & 52.0150 & 30.9776 & 2.73 & 28.0 & 59 & 59 \\
97 & 52.0284 & 29.4441 & 1.16 & -247.8 & 189 & 630 \\
98 & 52.0941 & 29.0017 & 2.62 & -194.4 & 52 & 181 \\
99 & 52.0991 & 29.0052 & 2.62 & -227.2 & 54 & 209 \\
100 & 52.1048 & 30.7135 & 2.84 & 25.2 & 53 & 52 \\
101 & 52.1113 & 29.0181 & 2.62 & -223.5 & 60 & 208 \\
102 & 52.1147 & 30.8918 & 3.81 & 87.2 & 50 & 47 \\
103 & 52.1211 & 28.9815 & 2.10 & -21.0 & 67 & 68 \\
104 & 52.1263 & 30.8459 & 5.16 & 47.9 & 36 & 31 \\
105 & 52.1315 & 30.8459 & 5.16 & 47.2 & 35 & 30 \\
106 & 52.1388 & 30.8443 & 5.16 & 47.7 & 36 & 31 \\
107 & 52.1503 & 29.3655 & 2.32 & -78.3 & 50 & 130 \\
108 & 52.1588 & 30.8360 & 4.21 & 50.6 & 35 & 32 \\
109 & 52.1682 & 30.8317 & 4.21 & 48.8 & 34 & 32 \\
110 & 52.1796 & 31.6535 & 4.46 & 184.6 & 216 & 46 \\
111 & 52.2177 & 29.5093 & 2.50 & -169.1 & 69 & 115 \\
112 & 52.2307 & 29.5287 & 2.90 & -124.0 & 46 & 203 \\
113 & 52.2943 & 31.5316 & 6.93 & 11.9 & 33 & 33 \\
114 & 52.2989 & 31.0010 & 3.37 & 67.5 & 57 & 55 \\
115 & 52.3346 & 31.1124 & 4.07 & 23.6 & 39 & 38 \\
116 & 52.3786 & 31.3691 & 4.97 & 140.4 & 61 & 50 \\
117 & 52.4963 & 30.6062 & 2.80 & -39.7 & 56 & 62 \\
118 & 52.4998 & 30.6105 & 2.64 & -31.3 & 55 & 58 \\
119 & 52.5153 & 30.5489 & 3.74 & -91.2 & 50 & 100 \\
120 & 52.5160 & 30.5547 & 3.74 & -104.6 & 52 & 113 \\
121 & 52.5203 & 31.1963 & 2.90 & -122.4 & 68 & 168 \\
122 & 52.5215 & 30.5514 & 3.74 & -45.4 & 38 & 58 \\
123 & 52.5238 & 30.6039 & 2.39 & -46.3 & 64 & 93 \\
124 & 52.5331 & 30.5480 & 3.75 & -25.9 & 34 & 42 \\
125 & 52.5379 & 30.5471 & 3.75 & -24.5 & 34 & 41 \\
126 & 52.5403 & 31.4094 & 2.05 & -310.9 & 121 & 87 \\
127 & 52.5764 & 30.9263 & 1.89 & -374.6 & 223 & 158 \\
128 & 52.5801 & 31.7726 & 2.77 & -199.5 & 86 & 283 \\
129 & 52.6025 & 31.7878 & 2.98 & -142.3 & 60 & 223 \\
130 & 52.6538 & 31.8737 & 2.56 & -276.4 & 103 & 501 \\
131 & 52.6616 & 31.5128 & 2.19 & -183.9 & 90 & 93 \\
132 & 52.6707 & 30.1919 & 0.79 & -416.7 & 550 & 431 \\
133 & 52.7403 & 31.5680 & 3.24 & -32.9 & 57 & 60 \\
134 & 52.7505 & 31.3681 & 1.93 & -148.8 & 95 & 413 \\
135 & 52.7948 & 31.6515 & 3.69 & -147.7 & 50 & 93 \\
136 & 52.7972 & 32.0584 & 1.49 & -197.8 & 156 & 139 \\
137 & 52.8051 & 31.9214 & 1.65 & 89.2 & 81 & 86 \\
138 & 52.8248 & 30.7921 & 2.85 & 82.7 & 82 & 54 \\
139 & 52.8296 & 30.7907 & 1.82 & 135.7 & 99 & 108 \\
140 & 52.8428 & 31.0572 & 1.95 & -109.9 & 90 & 319 \\
141 & 52.8888 & 31.8530 & 1.64 & 153.8 & 118 & 117 \\
142 & 52.9096 & 31.9073 & 1.99 & 21.9 & 76 & 74 \\
143 & 52.9134 & 30.6760 & 2.12 & -76.9 & 73 & 65 \\
144 & 52.9144 & 30.6719 & 2.12 & -66.9 & 66 & 65 \\
145 & 52.9186 & 30.6755 & 2.12 & -75.2 & 72 & 64 \\
146 & 52.9761 & 31.0781 & 2.20 & -301.5 & 148 & 127 \\
147 & 53.0330 & 30.4347 & 1.78 & 58.6 & 112 & 111 \\
148 & 53.0450 & 30.6622 & 2.26 & 55.2 & 79 & 75 \\
149 & 53.0664 & 31.7577 & 1.78 & 118.6 & 86 & 91 \\
150 & 53.0749 & 31.7560 & 1.69 & 108.0 & 100 & 89 \\
151 & 53.1052 & 31.0868 & 3.46 & 44.3 & 50 & 38 \\
152 & 53.1497 & 31.8378 & 1.73 & 352.6 & 581 & 117 \\
153 & 53.1670 & 31.8621 & 1.50 & 300.4 & 449 & 122 \\
154 & 53.1675 & 31.8530 & 1.50 & 310.5 & 464 & 124 \\
155 & 53.1680 & 31.8338 & 1.67 & 179.5 & 272 & 110 \\
156 & 53.1749 & 30.3356 & 2.27 & -38.3 & 53 & 54 \\
157 & 53.1795 & 31.8492 & 1.50 & 376.8 & 562 & 146 \\
158 & 53.1900 & 31.4131 & 3.90 & 34.6 & 52 & 50 \\
159 & 53.2240 & 31.7203 & 2.17 & 62.0 & 89 & 58 \\
160 & 53.2537 & 31.3384 & 5.32 & -104.4 & 46 & 50 \\
161 & 53.2750 & 30.8791 & 3.74 & 22.3 & 42 & 38 \\
162 & 53.2860 & 30.8468 & 2.22 & 205.5 & 95 & 123 \\
163 & 53.3239 & 30.7891 & 2.91 & -3.4 & 64 & 64 \\
164 & 53.3478 & 31.8360 & 1.06 & 566.2 & 1552 & 400 \\
165 & 53.3771 & 30.1060 & 1.31 & -31.0 & 122 & 164 \\
166 & 53.3779 & 31.2996 & 10.37 & -72.5 & 34 & 34 \\
167 & 53.4622 & 30.3840 & 1.89 & -68.5 & 98 & 467 \\
168 & 53.4651 & 30.5202 & 1.69 & 107.0 & 148 & 98 \\
169 & 53.4731 & 31.2204 & 4.28 & -16.6 & 46 & 48 \\
170 & 53.4811 & 29.8010 & 0.94 & -504.4 & 408 & 476 \\
171 & 53.4896 & 30.3951 & 1.53 & -43.4 & 112 & 238 \\
172 & 53.5077 & 31.3111 & 6.17 & -35.8 & 29 & 30 \\
173 & 53.5223 & 30.5026 & 1.14 & -57.2 & 162 & 163 \\
174 & 53.5459 & 30.3858 & 1.56 & -246.7 & 127 & 1244 \\
175 & 53.5491 & 31.0069 & 2.91 & -63.0 & 46 & 74 \\
176 & 53.5659 & 30.7473 & 1.31 & 489.9 & 1852 & 271 \\
177 & 53.5710 & 31.2029 & 4.09 & -55.1 & 32 & 43 \\
178 & 53.5778 & 30.8547 & 1.58 & -201.4 & 133 & 1199 \\
179 & 53.5802 & 30.8507 & 1.58 & -198.6 & 131 & 1182 \\
180 & 53.6689 & 30.6153 & 1.19 & 272.9 & 255 & 257 \\
181 & 53.7031 & 31.1146 & 2.54 & 36.1 & 53 & 54 \\
182 & 53.7377 & 30.8237 & 2.13 & 102.6 & 76 & 72 \\
183 & 53.7721 & 30.7852 & 2.04 & -84.2 & 67 & 68 \\
184 & 53.7787 & 30.7850 & 2.04 & -109.1 & 71 & 74 \\
187 & 53.8006 & 30.7847 & 1.70 & 61.8 & 78 & 82 \\
188 & 53.8651 & 30.3989 & 0.85 & 195.7 & 378 & 299 \\
189 & 53.8766 & 30.5400 & 2.11 & -90.7 & 58 & 104 \\
190 & 53.8765 & 30.7783 & 2.21 & 54.6 & 85 & 54 \\
191 & 53.8769 & 30.6246 & 1.55 & -134.4 & 90 & 91 \\
192 & 53.8881 & 31.5752 & 1.65 & 168.3 & 109 & 109 \\
193 & 53.8906 & 31.5713 & 1.65 & 145.8 & 103 & 103 \\
194 & 53.9287 & 30.3202 & 1.74 & -43.5 & 88 & 313 \\
195 & 54.2008 & 29.9774 & 1.21 & -201.9 & 161 & 756 \\
196 & 54.3165 & 31.2493 & 3.92 & 184.4 & 58 & 43 \\
200 & 54.4577 & 31.2542 & 2.42 & 358.6 & 52 & 110 \\
201 & 54.4661 & 30.9337 & 3.08 & 83.2 & 77 & 38 \\
202 & 54.5013 & 31.1029 & 2.45 & 158.4 & 113 & 59 \\
203 & 54.6136 & 30.8844 & 2.38 & 153.7 & 99 & 76 \\
204 & 54.9238 & 30.8482 & 1.71 & 301.2 & 225 & 145 \\
\end{longtable}

%##############################################
\section{Discussion}
\label{sec:discussion}

[TODO: Discussion section to be written - will include interpretation of field reversal, comparison with plane-of-sky field structure, implications for cloud formation scenarios, etc.]

%##############################################
\section{Conclusions and Future Work}
\label{sec:conclusions}

[TODO: Conclusions to be written]

%##############################################
\section{Software}
The \texttt{MC-BLOS} software is publicly available at: \url{https://github.com/MehrnooshTahani/MC-BLOS}. All data analysis and figure generation were performed using Python 3.9 with the following packages: NumPy \citep{harris2020array}, Pandas \citep{mckinney2010data}, Matplotlib \citep{hunter2007matplotlib}, and Astropy \citep{astropy2013, astropy2018}.

%##############################################
\section{Acknowledgment}

[TODO: Add acknowledgments]

%##############################################



\bibliography{main.bib}{}
\bibliographystyle{aasjournal}

\end{document}
